\chapter{Understanding the Source}
\label{chap:source-anatomy}
\index{command}\index{environment}\index{preamble}\index{argument!optional}\index{argument!required}

You can now compile a document. The next step is to read its source as a set of
small, recognisable structures. This skill is more valuable than memorising a
large preamble. When you can identify a command, its information, and the
region in which it acts, you can usually ask a precise question and test a
small repair.

We shall examine punctuation that \LaTeX{} treats specially. Nothing here
requires programming experience. The symbols are simply part of the written
grammar used to communicate with the compiler.

\begin{learningoutcomes}
  \item Distinguish commands, required arguments, and optional arguments.
  \item Match the beginning and end of an environment.
  \item Use comments to annotate source safely.
  \item Predict how spaces, line endings, and blank lines affect prose.
  \item Type common reserved characters in ordinary text.
  \item Separate the preamble from the document body.
  \item Construct a minimal working example and repair two common syntax
        errors.
\end{learningoutcomes}

\section{Commands: instructions beginning with a backslash}

In our first source, \cmd{documentclass}, \cmd{begin}, and \cmd{end} are
\term{commands}. A command tells \LaTeX{} to perform an operation or insert
something that ordinary text alone cannot express. Most command names made
from letters are case-sensitive. \cmd{begin} and \code{\textbackslash Begin}
are not the same instruction.

A command name made from letters ends when \LaTeX{} reaches a non-letter. In this
example, the opening brace ends the name \code{documentclass}:

\begin{latexcode}{A command and its information}
\documentclass{article}
\end{latexcode}

The backslash is not decoration. Without it, \LaTeX{} sees ordinary characters.
This is why removing the backslash in \cref{chap:first-document} changed the
meaning of the line.

\begin{classroomnote}
Not every backslash is followed by a long word. Some commands use a single
non-letter symbol. We will meet them when a task requires them. For now, read
``backslash plus a name'' as the usual command pattern.
\end{classroomnote}

\section{Required arguments use braces}

An \term{argument} supplies information to a command. A \term{required
argument} is normally enclosed in curly braces:

\begin{latexcode}{Required argument pattern}
\commandname{required information}
\end{latexcode}

The pattern line is explanatory rather than a standalone document. Our actual
class declaration has the same parts: \cmd{documentclass} is the command and
\code{article} is its required argument. Braces tell \LaTeX{} exactly which
characters belong to that argument.

Commands can have more than one argument, and an argument can contain other
commands. We will introduce those cases one at a time. The immediate habit is
to find each opening brace and its matching closing brace.

\begin{pausepredict}
Compare these two lines. Which one has a complete required argument?

\begin{latexcode}{One complete and one incomplete line}
\documentclass{article}
\documentclass{article
\end{latexcode}

The first has a matching pair of braces. The second leaves the argument open,
so \LaTeX{} may continue reading later lines while searching for the missing
closing brace.
\end{pausepredict}

\section{Optional arguments use brackets}

Some commands accept additional choices in square brackets. These are
\term{optional arguments}. The class command can accept options before its
required class name:

\begin{latexcode}{Complete example: an optional class argument}
\documentclass[11pt]{article}
\begin{document}
This document requests 11-point base text.
\end{document}
\end{latexcode}

Here \optional{11pt} is optional and \required{article} is required. If we
omit \optional{11pt}, the class uses its default base size. The square
brackets must appear in the position defined by the command. They are not a
general substitute for braces.

\expectedoutput

The PDF contains one sentence. Compared with the default minimal article, the
base type should be slightly larger. The class still controls the margins,
paragraph treatment, and page number.

\begin{tryit}
Compile the example with \optional{10pt}, \optional{11pt}, and
\optional{12pt}, one at a time. Keep all other source unchanged. Compare the
same sentence and the space it occupies. Which change is easiest to see?
\end{tryit}

\section{Environments have a beginning and an end}

An \term{environment} applies a structure or behaviour to a region. Its
general shape is:

\begin{latexcode}{Environment pattern}
\begin{environment-name}
Content inside the environment
\end{environment-name}
\end{latexcode}

The name after \cmd{begin} must match the name after \cmd{end}. Our document
body is the region inside the \code{document} environment. Later, other
environments will create lists, quotations, equations, tables, figures, and
proofs.

Environments can be nested, like boxes placed inside larger boxes. The inner
environment must end before the outer environment ends. Indentation is not a
\LaTeX{} requirement, but it helps human readers see that nesting.

\begin{commonerror}
\textbf{Broken example: mismatched names}

\begin{latexcode}{Broken example: do not compile as working source}
\documentclass{article}
\begin{document}
The names do not match.
\end{article}
\end{latexcode}

\textbf{Cause:} the source opens \code{document} but tries to close
\code{article}. A class name is not an environment name.

\textbf{Correction:} replace the final line with
\cmd{end}\required{document}. When the names match, compile again.
\end{commonerror}

\section{Comments explain source to people}

The percent sign begins a \term{comment}. \LaTeX{} ignores the percent sign and
the remainder of that source line:

\begin{latexcode}{Complete example: comments}
\documentclass{article} % Select the standard article class.
\begin{document}
This sentence appears in the PDF.
% This reminder does not appear in the PDF.
\end{document}
\end{latexcode}

Comments can explain a decision, mark a task, or temporarily remove a line
from a test. They should help a future reader understand \emph{why} the source
is organised as it is. A comment that merely repeats the command adds little.

Because the percent sign starts a comment, a literal percentage in prose needs
a backslash: \code{25\textbackslash\%} produces 25\%. We will return to
scientific percentages and quantities in the chapter on units.

\begin{goodpractice}
Write comments that survive beyond the current minute. ``Use the wide figure
because labels become unreadable at column width'' records a reason. ``Make
figure wide'' records only an instruction whose purpose may be forgotten.
Remove stale comments before submission.
\end{goodpractice}

\section{Whitespace: spaces, lines, and paragraphs}

In ordinary prose, \LaTeX{} usually treats one or several spaces as a single
interword space. A single source line ending also behaves like a space. A
blank line ends the current paragraph and begins another.

\begin{latexcode}{Complete example: source spacing}
\documentclass{article}
\begin{document}
Method A and     Method B
appear in one paragraph even though
the source changes line.

This sentence begins a new paragraph.
\end{document}
\end{latexcode}

\expectedoutput

The first three source lines form one typeset paragraph with ordinary word
spacing. The blank line creates a paragraph break before the final sentence.
\LaTeX{} does not print the row of extra source spaces.

This behaviour lets us wrap source at readable widths. Pressing Return is not
a request for a printed line break. In normal prose, use a blank line when you
mean a new paragraph. Commands for deliberate line breaks exist, but they
should solve a real structural need rather than imitate a typewriter.

\begin{pausepredict}
Remove the blank line before the final sentence and compile. Will there be one
paragraph or two? Now restore the blank line and add three more empty source
lines. Will the printed gap become three times larger? Test both predictions.
\end{pausepredict}

Ordinarily, several blank source lines still mean one paragraph break. If a
document needs more vertical space, its structure or design should specify the
reason. Repeated empty lines are not a reliable page-layout method.

\section{Characters with special jobs}

Several ASCII characters are \term{reserved characters}: \LaTeX{} uses them as
part of its instruction language. To print many of them literally, place a
backslash before the character.

\begin{table}[htbp]
  \centering
  \caption{Common reserved characters and a source form that prints them in
  ordinary text. The last three use named commands.}
  \label{tab:reserved-characters}
  \begin{tabular}{@{}lll@{}}
    \toprule
    Character & Common role in source & Source used to print it \\
    \midrule
    \# & numbered parameter in definitions & \code{\textbackslash\#} \\
    \$ & mathematics delimiter & \code{\textbackslash\$} \\
    \% & comment marker & \code{\textbackslash\%} \\
    \& & table alignment marker & \code{\textbackslash\&} \\
    \_ & subscript in mathematics & \code{\textbackslash\_} \\
    \{ and \} & argument grouping &
      \code{\textbackslash\{} and \code{\textbackslash\}} \\
    \textbackslash & command marker & \cmd{textbackslash} \\
    \textasciicircum & superscript marker & \cmd{textasciicircum} \\
    \textasciitilde & unbreakable space / active symbol &
      \cmd{textasciitilde} \\
    \bottomrule
  \end{tabular}
\end{table}

The special role can depend on context. For example, an ampersand will later
separate columns in a table or alignment. Outside such a structure, an
unescaped ampersand usually produces an error.

\begin{latexcode}{Complete example: literal reserved characters}
\documentclass{article}
\begin{document}
Recovery was 75\%, and A\&B shared sample\_01.

The budget was \$500 for item \#3.
\end{document}
\end{latexcode}

\expectedoutput

The PDF prints the percentage sign, ampersand, underscore, dollar sign, and
hash symbol. The backslashes used to escape them are not printed.

\begin{scienceinpractice}
Scientific source often contains reserved characters: \code{95\%} for a
percentage, \code{sample\_A1} for a file-like identifier, and \code{R\&D} in
an institutional phrase. Learning to recognise the character's source role
prevents apparently ordinary data labels from becoming syntax errors.
\end{scienceinpractice}

\section{The preamble and the body}

The \term{preamble} begins at the first source line and ends immediately before
\cmd{begin}\required{document}. It declares the class and later will load
packages, define reusable commands, and set document-wide behaviour. The
\term{document body} contains the material that forms the document's content.

\begin{latexcode}{Preamble and body}
\documentclass[11pt]{article} % Preamble

\begin{document}              % Body begins
This sentence is document content.
\end{document}                % Body ends
\end{latexcode}

The blank line in the tiny preamble is for human readability; it is not a
paragraph because typeset body content has not begun. Ordinary visible prose
does not belong in the preamble.

\begin{commonerror}
\textbf{Symptom:} ``Missing \cmd{begin}\required{document}'' appears even
though an opening line exists farther down.

\textbf{Possible cause:} ordinary text was accidentally placed before the
opening line. \LaTeX{} encountered content where it expected setup instructions.

\textbf{Repair:} move the prose into the body, then compile. If the message
remains, compare the spelling and braces of the opening line with the minimal
working document.
\end{commonerror}

\section{A minimal working example}

A \term{minimal working example}, often shortened to \term{MWE}, is the
smallest complete source that still demonstrates a behaviour or problem.
``Working'' means complete enough to compile until the target issue occurs;
it does not promise error-free output when the issue itself is an error.

An MWE helps in three ways:

\begin{itemize}
  \item it separates the problem from unrelated chapters and packages;
  \item it lets another person reproduce exactly what you observed; and
  \item it often reveals the cause while you are removing irrelevant source.
\end{itemize}

Suppose an ampersand causes an error in a 150-page report. A useful MWE might
be only this:

\begin{latexcode}{Broken MWE: unescaped ampersand}
\documentclass{article}
\begin{document}
Method A & Method B
\end{document}
\end{latexcode}

The example is complete, and the unescaped ampersand preserves the problem.
The repair is precise:

\begin{latexcode}{Corrected MWE}
\documentclass{article}
\begin{document}
Method A \& Method B
\end{document}
\end{latexcode}

Do not send an entire confidential thesis when five non-confidential lines can
demonstrate a technical fault. Replace sensitive content with a neutral example
and confirm that the fault remains.

\section{The first-aid debugging routine}

Let us apply a calm routine to a missing brace.

\begin{latexcode}{Broken example: missing closing brace}
\documentclass{article}
\begin{document}
This source has an unfinished group.
\end{document
\end{latexcode}

The compiler reaches the end while still looking for the closing brace of the
argument to \cmd{end}. It may report a runaway argument, an emergency stop, or
an unexpected end of file. The reported line can be at or after the true
cause, because \LaTeX{} noticed the problem only when it could no longer continue.

Use this sequence:

\begin{enumerate}
  \item Read the first meaningful error, not only the last line of the log.
  \item Note the reported line and inspect it plus a few preceding lines.
  \item Match braces and environment names.
  \item Make one small correction: add the closing brace.
  \item Save and compile again.
  \item Inspect the PDF; do not stop at a successful status message.
\end{enumerate}

\begin{latexcode}{Corrected closing line}
\end{document}
\end{latexcode}

If the small correction works, it provides evidence for the cause. If it does
not, return to the first new meaningful error. Avoid changing several unrelated
lines, because then you cannot tell which change mattered.

\begin{tryit}
Make a safe copy of your first document. Delete the closing brace from
\cmd{documentclass}\required{article}, compile, and record the first useful
message. Restore the brace and compile again. Compare this log with the
missing-brace example above: the wording may differ because the unfinished
argument begins in a different place.
\end{tryit}

\section{Reading source from left to right}

When a line looks unfamiliar, ask four questions:

\begin{enumerate}
  \item Where does the command name begin and end?
  \item Which braces contain required information?
  \item Are there square brackets containing an optional choice?
  \item Does the line open or close an environment that must be matched?
\end{enumerate}

Apply the questions to:

\begin{latexcode}{Source-reading practice}
\documentclass[12pt]{article}
\end{latexcode}

The command is \cmd{documentclass}. The optional argument is
\optional{12pt}. The required argument is \required{article}. The line does
not open an environment. This mechanical reading is simple, but it scales to
longer commands introduced later.

\begin{goodpractice}
Keep the source readable for a colleague: indent nested regions consistently,
place logical units on separate lines, and avoid dense sequences of commands
when a clearer arrangement is available. \LaTeX{} may ignore some whitespace;
human readers do not.
\end{goodpractice}

\chapterproject

\begin{miniproject}
Replace the running project's \file{main.tex} with the annotated version below.
The document still uses only structures already introduced.

\begin{latexcode}{Running project: readable source anatomy}
\documentclass[11pt]{article}

% Educational example: the values are simulated.
\begin{document}
This project compares Method A and Method B.
The eventual aim is a reproducible scholarly article.

All values are simulated demonstration data.
They do not describe a real experiment or clinical study.

Planned completeness is 100\%, and the project folder is
measurement\_methods.
\end{document}
\end{latexcode}

Before compiling, predict the number of paragraphs and identify every escaped
reserved character. Compile, inspect, and then write one helpful comment of
your own. Update the README only if your new comment records information that
belongs at project level rather than beside one source decision.
\end{miniproject}

\chaptersummary

Commands usually begin with a backslash. Braces group required arguments;
square brackets commonly contain optional arguments. Environments have
matching beginning and ending names. A percent sign starts a source comment,
while reserved characters need special treatment when printed literally.
Ordinary spaces and line endings collapse in prose, but a blank line starts a
new paragraph. The preamble configures the document; the body contains its
content. A minimal working example and the first-meaningful-error routine turn
a large mystery into a small test.

\chaptercheck

\begin{chapterchecklist}
  \item I can identify a command name and its arguments.
  \item I match every opening environment with the same closing name.
  \item I can add a comment and explain why it does not print.
  \item I can predict whether source lines form one paragraph or two.
  \item I can print common reserved characters safely.
  \item I can point to the preamble and body of a complete source file.
  \item I can reduce a simple error to a minimal working example.
\end{chapterchecklist}

\chapterexercisestitle

\begin{chapterexercises}
  \item Identify the command, optional argument, and required argument in
        \code{\textbackslash documentclass[12pt]\{book\}}.
  \item Explain why braces and brackets are not interchangeable in the class
        declaration.
  \item Predict the PDF paragraph count when five source lines contain no
        blank line. Test your answer.
  \item Write source that prints: \code{Accuracy was 95\% for sample\_7.}
  \item Find and repair the mismatch:
        \code{\textbackslash begin\{document\}} followed by
        \code{\textbackslash end\{article\}}.
  \item Create a five-line MWE that demonstrates an unescaped ampersand error,
        then create the corrected version.
  \item Read this error strategy: ``Change five nearby lines and compile once.''
        Explain why it produces weak evidence about the cause.
  \item Challenge: take a harmless problem from your practice source and
        reduce it until removing any remaining line would stop the example
        from demonstrating the same problem.
\end{chapterexercises}
